\chapter{运动方程与守恒律}

\section{从物理模型到微分方程}
把物体视为质点，意味着它的大小和内部运动不影响当前关心的平移运动。
这并不要求物体真的没有尺寸：计算行星轨道时可以忽略自转，计算陀螺时就不能这样做。
在惯性系中，对质量恒定的质点，
\begin{equation}
 m\ddot{\bm r}=\bm F(\bm r,\dotbm{r},t).
 \label{eq:cm-newton}
\end{equation}
给定初始位置和速度，便得到一个初值问题。在通常的光滑力模型下，运动在局部唯一。\footnote{一组充分条件是力连续，且在所考察的时间邻域内对$(\bm r,\dotbm{r})$一致局部Lipschitz。}
碰撞、奇异势和分段摩擦需要另加处理，不能对任意方程都预先断言解存在且唯一。

实际解题的顺序是：选系统和参考系，画受力图，写约束，再选择坐标。
支持力和张力来自约束，并不总能事先取成$mg$；静摩擦力满足不等式$\abs{f}\leq\mu_s N$，
只有临界滑动时才取等号。滑动摩擦$\mu_kN$则是经验模型，方向与接触面的相对滑动速度相反。

两个惯性系以恒定相对速度$\bm u$平移时，$t'=t$，$\bm r'=\bm r-\bm u t-\bm r_0$，
因此$\ddot{\bm r}'=\ddot{\bm r}$。这是Galileo变换下Newton方程的运动学基础。
速度、动能依赖参考系；力学规律的形式相同不意味着每个物理量的数值相同。

\section{平面极坐标中的速度与加速度}
笛卡尔坐标中的基矢固定，故可以直接逐分量求导。平面极坐标的基矢随粒子位置旋转：
\begin{align}
 \bm r&=r\uvec r,&
 \dot{\uvec r}&=\dot\phi\uvec\phi,&
 \dot{\uvec\phi}&=-\dot\phi\uvec r,\\
 \bm v&=\dot r\uvec r+r\dot\phi\uvec\phi,&&&
 \bm a&=(\ddot r-r\dot\phi^2)\uvec r
 +(r\ddot\phi+2\dot r\dot\phi)\uvec\phi.
 \label{eq:cm-polar-a}
\end{align}
基矢导数可由$\uvec r=(\cos\phi,\sin\phi)$直接得到；向心项不是额外施加的力。
对半径固定的圆周运动，径向加速度为$-r\dot\phi^2$，切向加速度为$r\ddot\phi$。
以后选择极坐标，是因为中心力天然区分径向与切向，而不是为了增加坐标形式。

\section{功、势能与能量守恒}
将式\eqref{eq:cm-newton}与速度点乘，得到功率关系
\begin{equation}
 \dd{t}\Bigp{\frac12m v^2}=\bm F\cdot\bm v,
 \qquad T_B-T_A=\int_A^B\bm F\cdot\d\bm r.
\end{equation}
若$\bm F=-\grad V(\bm r)$，则$E=T+V$守恒。
在单连通区域内，连续可微力场的$\curl\bm F=0$保证存在单值势能；
在有洞区域内还须检查闭合回路的环量，不能仅凭局部旋度判断全局保守性。

更一般地，若$\bm F=-\grad V(\bm r,t)+\bm F_\nonconservative$，则
\begin{equation}
 \dot E=\pv{V}{t}+\bm F_\nonconservative\cdot\bm v.
 \label{eq:cm-energy-work}
\end{equation}
势能显含时间时，即使力在每一时刻都可由势函数给出，机械能也未必守恒。
磁力$q\bm v\times\bm B$不做功，却通常不能用仅依赖位置的标量势表示；“不做功”和“保守力”不是同一个条件。

一维保守运动满足$\dot x=\pm\sqrt{2[E-V(x)]/m}$。
允许区域由$V(x)\leq E$给出，转向点满足$V(x)=E$。
平衡点满足$V'(x_0)=0$；若$V''(x_0)>0$，线性化给出
\begin{equation}
 m\ddot\eta+V''(x_0)\eta=0,\qquad
 \omega_0^2=V''(x_0)/m,
 \quad \eta=x-x_0.
\end{equation}
若二阶导数为零，需要保留高阶项。例如$V=\alpha x^4$在$\alpha>0$时仍稳定，
但没有非零的小振幅线性频率。

\begin{example}{线性阻力下的落体}{cm-drag}
取竖直向下为正，阻力$-bv$，$b>0$。由$m\dot v=mg-bv$得
\begin{equation}
 v(t)=v_\infty+(v_0-v_\infty)\e^{-t/\tau},\quad
 v_\infty=\frac{mg}{b},\quad \tau=\frac mb.
\end{equation}
积分一次给出$x(t)-x_0=v_\infty t+(v_0-v_\infty)\tau(1-\e^{-t/\tau})$。
从静止开始时，$t\ll\tau$恢复$x-x_0\simeq gt^2/2$，$t\gg\tau$趋于匀速。
这里线性阻力是适用速度范围内的模型；高速空气阻力常需改为近似正比于$v^2$。
\end{example}

\section{质点组、质心与角动量}
对固定成员、质量不变的$N$个粒子，令
\begin{equation}
 M=\sum_a m_a,\quad \bm R=\frac1M\sum_a m_a\bm r_a,
 \quad \bm P=\sum_a m_a\dotbm{r}_a=M\dotbm{R}.
\end{equation}
若内部力成对大小相等、方向相反，求和得到
\begin{equation}
 \dotbm{P}=\bm F_\external,\qquad M\ddot{\bm R}=\bm F_\external.
\end{equation}
因此无外力时质心匀速运动。令$\bm\rho_a=\bm r_a-\bm R$，则
\begin{equation}
 T=\frac12M\norm{\dotbm{R}}^2+\frac12\sum_a m_a\norm{\dotbm{\rho}_a}^2.
 \label{eq:cm-com-energy}
\end{equation}
交叉项因$\sum_a m_a\bm\rho_a=0$而消失。这区分整体平移与内部运动。

关于惯性系中固定原点的角动量为$\bm\ell=\sum_a\bm r_a\times m_a\dotbm{r}_a$。
若内部力还沿两粒子连线，则其内力矩相消，
\begin{equation}
 \dotbm{\ell}=\bm\tau_\external,\qquad
 \bm\ell=\bm R\times M\dotbm{R}+\bm\ell_\CM.
\end{equation}
第一项是质心的轨道角动量，第二项是关于质心的角动量；
关于质心同样有$\dotbm{\ell}_\CM=\bm\tau_{\external,\CM}$。
对任意运动的取矩点则须重新求导，不能直接套用固定原点公式。
含电磁场的系统还可能需要计入场动量，不能只对粒子强行假定内力遵循上述成对关系。

\section{碰撞和变质量：先选闭合系统}
短时间碰撞中，若外力冲量可忽略，则总动量守恒；机械能只在弹性碰撞中守恒。
一维两物体碰撞可用恢复系数$e$补充关系
\begin{equation}
 v_2'-v_1'=-e(v_2-v_1),\qquad
 \Delta T=-\frac12\mu(1-e^2)(v_1-v_2)^2,
 \quad \mu=\frac{m_1m_2}{m_1+m_2}.
\end{equation}
这里碰前两物体相互接近，$0\leq e\leq1$是常见被动碰撞模型；$e=0$时碰后两者速度相同。
质心动能不变，损失来自相对运动的动能。

火箭不是固定成员的粒子集合。取前进方向为正，设喷气相对火箭的速度为$-u$，$u>0$，
在$\d t$内排出质量$-\d m>0$。对“火箭加刚排出的气体”用动量定理，保留一阶小量：
\begin{equation}
 mv+F_\external\d t=(m+\d m)(v+\d v)+(-\d m)(v-u),
\end{equation}
于是
\begin{equation}
 m\dot v=F_\external-u\dot m.
\end{equation}
无外力且$u$恒定时，$v_f-v_i=u\ln(m_i/m_f)$。
直接对火箭写$F_\external=\d(mv)/\nd t$会漏掉穿越系统边界的动量流。

\paragraph{复习与阅读}
Newton方程给出各分量的运动；守恒律减少要求解的自由度。
先检查系统边界和参考系，再判断能否使用动量、角动量、能量。
延伸阅读：Taylor的动量与角动量、能量章节；Morin的守恒律与变质量问题。
